\documentclass[../../main.tex]{subfiles}
\begin{document}
\section{函数的微分}\label{section:ch11-函数的微分}

在初等微积分中，$\mathbb{R}^n$ 上光滑实值函数 $f$ 的梯度定义为以 $f$ 的偏导数为分量的向量场。用我们的记号，它读作
\begin{align}
\mathrm{grad}\,f=\sum_{i=1}^n\frac{\partial f}{\partial x^i}\frac{\partial}{\partial x^i}.\label{align:ch11-梯度}
\end{align}
遗憾的是，这种形式的梯度不能与坐标无关地定义。（它违反我们指标约定这一事实是一条强有力的线索。）

\begin{exercise}\label{exercise:ch11-17}
设 $f(x,y)=x^2$ 在 $\mathbb{R}^2$ 上，$X$ 为向量场
\begin{align*}
X=\mathrm{grad}\,f=2x\frac{\partial}{\partial x}.
\end{align*}
在极坐标下（在它们有定义的某个开子集上）利用\eqref{align:ch11-余向量坐标变换}~计算 $X$ 的坐标表达式，并证明它不等于
\begin{align*}
\frac{\partial f}{\partial r}\frac{\partial}{\partial r}+\frac{\partial f}{\partial\theta}\frac{\partial}{\partial\theta}.
\end{align*}
\end{exercise}

虽然光滑函数的偏导数不能与坐标无关地解释为向量场的分量，但结果表明，它们可以非常自然地解释为余向量。下面给出精确的定义。

设 $M$ 为光滑流形，$f\in C^\infty(M)$。对每个 $p\in M$，我们定义一个线性映射 $df_p\colon T_pM\to\mathbb{R}$，称为 $f$ 在 $p$ 处的\textbf{微分}，定义为
\begin{align}
df_p(v)=v(f),\quad v\in T_pM.\label{align:ch11-微分定义}
\end{align}
由于它是 $T_pM$ 上的线性泛函，$df_p$ 是 $p$ 处的一个余向量。当 $p$ 变动时，这就定义了 $M$ 上的一个余向量场 $df$。我们称 $df$ 为 $f$ 的\textbf{微分}。下面的命题总结了它的性质。

\begin{proposition}\label{proposition:ch11-微分性质}
设 $M$ 为带边或不带边的光滑流形。微分映射 $d\colon C^\infty(M)\to\mathfrak X^*(M)$ 满足下列性质：
\begin{enumerate}[(a)]
\item $d$ 是线性的：$d(af+bg)=a\,df+b\,dg$；
\item 乘积法则：$d(fg)=f\,dg+g\,df$；
\item 对任意 $p\in M$，$df_p$ 是由\eqref{align:ch11-微分定义}~定义的余向量。
\end{enumerate}
\end{proposition}

\begin{proof}
线性与乘积法则都来自切向量作用于函数时的相应性质。
\end{proof}

在坐标图 $(U,\varphi)$ 中，坐标函数为 $(x^i)$，我们可计算 $df$ 的分量。若 $v=v^j\partial/\partial x^j|_p\in T_pM$，则
\begin{align*}
df_p(v)=v(f)=v^j\frac{\partial f}{\partial x^j}(p).
\end{align*}
另一方面，若把 $df$ 写为坐标余基 $(\theta^i)$ 的线性组合 $df=df_j\theta^j$，则
\begin{align*}
df_p(v)=df_j(p)\theta^j(v)=df_j(p)v^j.
\end{align*}
比较两式，可见
\begin{align}
df_j(p)=\frac{\partial f}{\partial x^j}(p).\label{align:ch11-微分坐标}
\end{align}
于是
\begin{align}
df=\frac{\partial f}{\partial x^i}\theta^i.\label{align:ch11-微分坐标形式}
\end{align}

特别地，把此公式应用于坐标函数 $x^i$ 本身，得到
\begin{align*}
dx^i=\frac{\partial x^i}{\partial x^j}\theta^j=\delta^i_j\theta^j=\theta^i.
\end{align*}
也就是说，坐标余向量场 $(\theta^i)$ 恰为坐标函数的微分。出于这个原因，从现在起我们弃用记号 $\theta^i$ 而改用 $dx^i$ 表示坐标余基。在这个新记号下，\eqref{align:ch11-微分坐标形式}~变为
\begin{align}
df=\frac{\partial f}{\partial x^i}dx^i.\label{align:ch11-微分df}
\end{align}

在 $\mathbb{R}^n$ 上，把微分解释为对 $\Delta f$ 的最佳线性近似是很自然的。若 $\Delta f=f(p+v)-f(p)$（对 $v\in\mathbb{R}^n$），Taylor 定理表明，当 $v$ 小时 $\Delta f$ 被很好地近似为
\begin{align*}
\Delta f\approx\frac{\partial f}{\partial x^i}(p)v^i=\frac{\partial f}{\partial x^i}(p)\,dx^i|_p(v)=df_p(v)
\end{align*}
（这里我们通过通常的等同把 $v$ 视为 $T_p\mathbb{R}^n$ 的元素）。换言之，$df_p$ 是在 $p$ 附近最佳近似 $\Delta f$ 的线性泛函（\cref{figure:ch11-微分近似}）。微分概念的巨大威力来自这样一个事实：我们可以在任意流形上不变地定义 $df$，而不必求助于涉及无穷小的模糊论证。

\begin{figure}[htbp]
\centering
\includegraphics[width=0.5\textwidth]{fig-11-2.png}
\caption{微分作为 $\Delta f$ 的近似}\label{figure:ch11-微分近似}
\end{figure}

下面的命题是微分版的\cref{proposition:ch03-24}。

\begin{proposition}[函数沿曲线的导数]\label{proposition:ch11-沿曲线导数}
设 $M$ 为带边或不带边的光滑流形，$\gamma\colon J\to M$ 是光滑曲线，$f\colon M\to\mathbb{R}$ 是光滑函数。则实值函数 $f\circ\gamma\colon J\to\mathbb{R}$ 的导数为
\begin{align}
(f\circ\gamma)'(t)=df_{\gamma(t)}(\gamma'(t)).\label{align:ch11-沿曲线导数}
\end{align}
\end{proposition}

\begin{proof}
见\cref{figure:ch11-沿曲线导数}。直接由定义，对任意 $t_0\in J$，
\begin{align*}
df_{\gamma(t_0)}(\gamma'(t_0))&=\gamma'(t_0)(f) &&\text{（$df$ 的定义）}\\
&=d\gamma_{t_0}\left(\frac{d}{dt}\bigg|_{t_0}\right)f &&\text{（$\gamma'(t)$ 的定义）}\\
&=\frac{d}{dt}\bigg|_{t_0}(f\circ\gamma) &&\text{（$d\gamma$ 的定义）}\\
&=(f\circ\gamma)'(t_0) &&\text{（$d/dt|_{t_0}$ 的定义）}.
\end{align*}
\end{proof}

\begin{figure}[htbp]
\centering
\includegraphics[width=0.7\textwidth]{fig-11-3.png}
\caption{函数沿曲线的导数}\label{figure:ch11-沿曲线导数}
\end{figure}

\begin{exercise}\label{exercise:ch11-24}
对光滑实值函数 $f\colon M\to\mathbb{R}$，证明 $p\in M$ 是 $f$ 的临界点当且仅当 $df_p=0$。
\end{exercise}

你可能已经注意到，对光滑实值函数 $f\colon M\to\mathbb{R}$，我们现在对 $f$ 在点 $p\in M$ 处的微分有了两种不同定义。在第~\cref{chapter:ch03}~中，我们把 $df_p$ 定义为从 $T_pM$ 到 $T_{f(p)}\mathbb{R}$ 的线性映射。在本章中，我们把 $df_p$ 定义为 $p$ 处的余向量，即从 $T_pM$ 到 $\mathbb{R}$ 的线性映射。一旦考虑 $\mathbb{R}$ 与 $T_{f(p)}\mathbb{R}$ 之间的典范等同，这两者实际上是同一个对象；一种容易看出这一点的方式是注意二者在坐标下都由 $f$ 的偏导数作为行矩阵表示。

类似地，若 $\gamma$ 是 $M$ 中的光滑曲线，表达式 $(f\circ\gamma)'(t)$ 对我们有两种不同含义。一方面，$f\circ\gamma$ 可解释为 $\mathbb{R}$ 中的光滑曲线，因而 $(f\circ\gamma)'(t)$ 是它在点 $f\circ\gamma(t)$ 处的速度，即切空间 $T_{f\circ\gamma(t)}\mathbb{R}$ 的元素。\cref{proposition:ch03-24}~表明此切向量等于 $df_{\gamma(t)}(\gamma'(t))$，视为 $T_{f\circ\gamma(t)}\mathbb{R}$ 的元素。另一方面，$f\circ\gamma$ 也可简单地视为一元实变量的实值函数，此时 $(f\circ\gamma)'(t)$ 就是它的普通导数。\cref{proposition:ch11-沿曲线导数}~表明此导数等于 $df_{\gamma(t)}(\gamma'(t))$，视为实数。我们选哪种解释取决于心中的目的。
\end{document}
